\documentclass[11pt,a4paper]{article}
\usepackage[T1]{fontenc}
\usepackage{lmodern,amsmath,amssymb,amsthm,mathtools,booktabs,microtype,xcolor}
\usepackage[margin=26mm]{geometry}
\usepackage[colorlinks=true,linkcolor=blue!45!black,citecolor=blue!45!black,urlcolor=blue!45!black]{hyperref}
\newtheorem{theorem}{Theorem}[section]
\newtheorem{lemma}[theorem]{Lemma}
\newtheorem{proposition}[theorem]{Proposition}
\newtheorem{corollary}[theorem]{Corollary}
\theoremstyle{remark}\newtheorem{remark}[theorem]{Remark}
\newcommand{\Q}{\mathbb Q}\newcommand{\Z}{\mathbb Z}
\DeclareMathOperator{\Tr}{Tr}\DeclareMathOperator{\N}{N}
\setlength{\emergencystretch}{2em}
\title{A quartic in 389 variables from local nonvanishing}
\author{}\date{8 November 2026}
\begin{document}\maketitle
\begin{abstract}
We specify an integer quartic in 389 total variables, with 6198 nonzero
monomials, by 361 quadratic equations and a complete coefficient table.
Its rational input projection is the integers under the same arithmetic
premises as the specified 392-variable construction. Three rational
inverse witnesses can be omitted: at every prime where a local
divisibility conclusion is required, the remaining equations already
force their denominators to be nonzero. The additional fibers have
empty contact radicals, for which the required divisibility statement
is vacuous. The argument explicitly covers the degenerate value at
which the detector's rational formula has a zero denominator.
The polynomial has content one and coefficient height less than
$2^{1791357}$. The largest coefficient is exactly three smaller than
that of the starting polynomial.
\end{abstract}

\section{The specified polynomial and its arithmetic scope}
For $F\in\Z[X,Y_1,\ldots,Y_m]$, its rational input projection is
\[
 \{x\in\Q:\exists\boldsymbol y\in\Q^m,\ F(x,\boldsymbol y)=0\}.
\]
Every variable count includes $X$. Degree means total degree, and
coefficient height means the largest absolute integer coefficient after
the fixed constants have been substituted.

\begin{theorem}\label{thm:main}
The files \nolinkurl{quadratic-system389.json} and
\nolinkurl{polynomial389.json} specify 361 integer quadratic polynomials
$Q_i$ in 389 coordinates and their sum of squares
$P_{389}=\sum_iQ_i^2$. Its degree is exactly four, its integer content
is one, and it has 6198 nonzero coordinate monomials. Its largest
coefficient has 1791357 binary digits.

Under the arithmetic premises of the 392-variable construction
\cite{Construction392}, its rational input projection is exactly $\Z$.
No additional prime-pattern or arithmetic existence assertion is used.
\end{theorem}

The retained arithmetic premises are the specified elliptic-group and
index-height assertions, the contact-height theorem, the pole obstruction
and simultaneous pole realization, and Sun's simultaneous prime-pattern
theorem with the quantitative weighted selection premise used in the
quartic construction. These hypotheses are stated precisely in
\cite[Section~2]{Construction392}. The present theorem proves a new
implication from those premises. The coefficient checks do not establish
the deep arithmetic inputs.

The constant in the coefficient table is
\[
                         A=80989^{54929}.
\]
It is an integer coefficient, not an additional coordinate. The starting
polynomial is fixed by commit
\begin{center}\small
\nolinkurl{de7aa60dd1794d7ea62a8a83cdc4bfd22380c74f}
\end{center}
of the repository in \cite{Construction392}. Its full quadratic system
and coefficient table are included as exact compressed source files.

\section{The detector equations}
Put $K=\Q(\rho)$, where $\rho^2=2$, and let $\sigma$ denote its
nontrivial automorphism. For $W=w_0+w_1\rho$ one has
\[
                       \N(W)=w_0^2-2w_1^2.
\]
In particular, $\N(W)=0$ if and only if $W=0$.

For the three retained height detectors write
\[
 a\in\{1,C_0^2,(1+C_1^2)^2\},\qquad t=a\lambda,
 \qquad\delta=t-1,
\]
where $C_0+C_1\rho$ is the product of the four root abscissae.
The retained root charts and their real embeddings give $C_0>0$;
all three values of $a$ are therefore nonzero rational squares.
The stored coordinates use affine substitutions. To remove any
ambiguity, if $u$ and $v$ are the coordinates named
\nolinkurl{heightProduct.Z.0} and \nolinkurl{heightProduct.Z.1}, then
\[
 C_0=2u,\qquad C_1=4v.
\]
The stored quantities named \nolinkurl{scale0},
\nolinkurl{x1Square}, and \nolinkurl{scale1} are $u^2,v^2,v^4$.
Their consumers, rather than these names alone, recover the actual
scales. Specifically, the stored coordinates $t_1^*,t_2^*$ satisfy
\[
 t_1^*=\lambda u^2,\qquad
 t_2^*=\lambda v^2+8\lambda v^4,
\]
and the actual inputs are
\[
 t_0=\lambda,\qquad t_1=4t_1^*=C_0^2\lambda,
 \qquad t_2=\lambda+32t_2^*=(1+C_1^2)^2\lambda.
\]
These identities follow by direct substitution into the retained
quadratic rows. The independent normalization checker verifies them
as polynomial identities, together with the weight and pole-input
identities below.
Set
\[
 \Theta=\left\{0,-\tfrac12,-1,\tfrac1{2399},-\tfrac{32}{407}\right\},
 \qquad F(t)=\prod_{\theta\in\Theta}(t-\theta).
\]
Let $M(t)$ be the product of the primes $p>76176$ for which
$v_p(t-\theta)$ is positive and odd for at least one
$\theta\in\Theta$.

With these actual coordinates, the retained quadratic gates imply
the polynomial identity
\begin{equation}\label{eq:e}
                            F(t)e=(a\delta)^5.
\end{equation}
Whenever $\delta F(t)\ne0$, put
\[
 R=\frac{F(t)}{a^5\delta^6},\qquad
 q_m=\frac{R^{1-2m}}{\delta^{2m+2}},
 \qquad
 W=\prod_j(1+q_{m_j}D_j^2)=w_0+w_1\rho.
\]
Here the $D_j\in K$ are the retained field-valued point, comparison,
input, and Lucas tests. Their assigned orders are $m_j\in\{1,10,40\}$.
The grouped products and power gates of the starting system compute
exactly this product. Their definitions are polynomial identities;
no deleted inverse witness is used by these gates.

The final detector rows, apart from its pole certificate, are
\begin{equation}\label{eq:trace-clearing}
 w_0^2-2w_0T-2w_1^2=0,
 \qquad e\delta b=T^2.
\end{equation}
The starting system additionally imposed $w_0I=1$ with a private
rational inverse witness $I$. We delete this last row and this
coordinate in each detector. The five normalized pole rows are
retained. Whenever $b\ne0$, they imply the same obstruction as the
original pole equations: $v_p(b)$ cannot be negative and odd for
$p\notin\{2,3,17\}$.

\section{Local nonvanishing without inverse witnesses}

\begin{lemma}\label{lem:contacts}
For every rational solution of the modified system, $F(t)\ne0$.
If $p\mid M(t)$, then $\delta\ne0$, $v_p(\delta)=0$, and
$v_p(R)$ is odd. If $\delta=0$, then $M(t)=1$.
\end{lemma}
\begin{proof}
If $F(t)=0$, then $t\in\Theta$, and hence $\delta\ne0$ because
$1\notin\Theta$. This contradicts \eqref{eq:e}, since $a\ne0$.

For $p>76176$, every nonzero difference of two elements of
$\Theta\cup\{1\}$ is a $p$-adic unit. This follows by direct
factorization of their rational numerators and denominators. Thus a
prime in $M(t)$ divides exactly one factor of $F(t)$ to positive odd
order, while $\delta$ is a unit. Since $a$ is a rational square,
\[
                    v_p(R)=v_p(F(t))-5v_p(a)
\]
is odd. Its sign has not yet been assumed.

Finally $\delta=0$ gives $t=1$. The five differences $1-\theta$ have
no numerator or denominator prime greater than 76176. Consequently
$M(1)=1$.
\end{proof}

\begin{lemma}\label{lem:local-nonzero}
If $M(t)>1$, then $W,w_0,T,b$ are nonzero and the retained detector
equations have their intended rational interpretation. In particular,
\[
              T=\frac1{\Tr(W^{-1})},\qquad
              b=\frac{R}{\Tr(W^{-1})^2}.
\]
\end{lemma}
\begin{proof}
Choose $p\mid M(t)$ and put $\nu=v_p(R)$, an odd integer by
Lemma~\ref{lem:contacts}. The prime $p$ is unramified in $K$. At every
place $v$ above $p$, normalized by $v(p)=1$, its value group on $K^\times$
is $\Z$. Since $\delta$ is a unit,
\[
             v(q_mD^2)=(1-2m)\nu+2v(D)
\]
is odd whenever $D\ne0$. It cannot be the valuation of $-1$.
Therefore every factor $1+q_mD^2$ is nonzero, also when $D=0$,
and hence $W\ne0$.

If $w_0=0$, the first equation of \eqref{eq:trace-clearing} gives
$w_1=0$, contrary to $W\ne0$. Thus $w_0\ne0$. Its norm is nonzero,
and the same equation gives
\[
 T=\frac{\N(W)}{2w_0}
   =\frac1{\Tr(W^{-1})}\ne0.
\]
Moreover \eqref{eq:e} gives $e\delta=1/R\ne0$. The second equation
of \eqref{eq:trace-clearing} now yields $b=RT^2\ne0$.
All divisions are justified by the remaining equations and the chosen
contact prime; none uses the deleted inverse row.
\end{proof}

\begin{proposition}\label{prop:local-tests}
For every $p\mid M(t)$ there is a place $v$ above $p$ at which
\[
                         v(D_j)\ge m_j\quad\text{for every }j.
\]
The asserted divisibility consequences of the starting detector remain
valid when $M(t)=1$.
\end{proposition}
\begin{proof}
By Lemma~\ref{lem:local-nonzero} the pole input and the norm-trace
formula are defined and nonzero. For each $v\mid p$, parity gives
\[
 v(1+q_{m_j}D_j^2)
 =\min\{0,\nu+2(v(D_j)-m_j\nu)\}\le0.
\]
Put $\ell=\max_{v\mid p}v(W)\le0$. The trace inequality implies
\[
 v_p(\Tr(W^{-1}))\ge-\ell,
 \qquad v_p(T)\le\ell.
\]
The valuation of $b=RT^2$ is odd. The pole obstruction therefore
implies
\[
                      0<v_p(b)\le\nu+2\ell.
\]
In particular $\nu>0$ and $\ell>-\nu/2$. At a place attaining
$\ell$, every factor valuation is at least $\ell$, because all the
factor valuations are nonpositive and their sum is $\ell$.
Consequently
\[
              v(D_j)>(m_j-3/4)\nu\ge m_j-3/4.
\]
The value group is integral, so $v(D_j)\ge m_j$.

If $M(t)=1$, every conclusion of the form $M(t)^k\mid N$ is
automatic. No application of a rational pole formula is needed for
that detector in this case. This includes $t=1$, for which $h(t)=0$.
\end{proof}

\begin{remark}
The argument is local in its use of nonvanishing but proves a global
integer-projection statement. It does not assert that every new rational
witness tuple lifts to one of the old tuples: a detector with empty
contact radical can have $W=0$. Such a tuple is handled by the vacuous
divisibility conclusion, not by an undefined reciprocal.
\end{remark}

\section{Soundness and completeness}

\begin{proof}[Proof of the projection assertion in Theorem~\ref{thm:main}]
All retained point charts, scalar equations, the forced-prime
normalizer, and the scalar nonzero guard are unchanged. The three
deleted coordinates occur only in their respective inverse rows.
The arithmetic point interpretations and the fixed-prime index estimate
therefore remain valid before any use of a height detector.

Proposition~\ref{prop:local-tests} gives exactly the local test orders
required at every prime in each of the three contact radicals. The
formal logarithm comparison, integer-label congruences and the two
height contradictions in \cite[Sections~2 and~5]{Construction392}
use these conclusions. They do not require a detector denominator at
primes outside its contact radical. They are consequently unchanged.
The contact-height bound itself is a statement about the rational
number $t$; it applies also when its contact radical is one. By
Lemma~\ref{lem:contacts}, none of its excluded contact values occurs.

In particular the unchanged argument proves that the Lucas residual
is zero and then that the input label
\[
                         \beta=51(2X+1)^2+2504
\]
equals its odd integer elliptic index. If $2X+1=u/v$ is in lowest
terms with $v>0$, this forces $v^2\mid51$, so $v=1$; oddness of
$\beta$ then makes $u$ odd. Hence $X\in\Z$.

Conversely every integer input has a solution of the 392-variable
system under the same arithmetic premises. Forgetting the three
private inverse coordinates gives a solution of the new system.
There are no new constraints and no new simultaneous-choice problem.
Finally a sum of rational squares is zero exactly when every square
is zero, proving the assertion for $P_{389}$.
\end{proof}

\section{Exact data and independent verification}
The source coordinate numbers and residual numbers are inherited from
the upstream arithmetic circuit. Their deletion is the complete change:
\begin{center}
\begin{tabular}{ccc}
\toprule Source coordinate & Coordinate in the 392 system & Source row\\
\midrule382&221&334\\430&250&384\\478&279&434\\\bottomrule
\end{tabular}
\end{center}
Each deleted row is $1-w_0I=0$. Every other primitive integer residual
is retained verbatim, apart from renumbering its coordinates in increasing
order. Thus the counts are $392-3=389$ coordinates and $364-3=361$
quadratic residuals.

The producer \nolinkurl{build_quartic389.py} performs this deletion
and expands the complete sum of squares. The separate program
\nolinkurl{check_quartic389.py} imports no producer. It reconstructs
the original 392-variable polynomial, checks the sole-use inverse
coordinates and every retained row, and compares the entire new
coefficient table in two ways: direct squaring and subtraction of
the three deleted guard squares.

The additional program \nolinkurl{verify_detector_normalizations.py}
independently substitutes the complete affine coordinate dictionary
into the relevant residuals, using exact rational arithmetic. The
three weight rows become $F(t)e-(a\delta)^5$ multiplied, respectively,
by $1952786$, $976393/8$, and $1952786$. The retained rows also give
$\delta^3q_1=e$ and $e\delta b=T^2$. The toric power equations give
$q_{10}=e^{19}/\delta^3$ and $q_{40}=e^{79}/\delta^3$ on their
nonzero domain. Thus the affine substitutions do not change the
semantic scales, weights, or pole inputs.

There are 6219 terms when powers of the fixed integer $A$ are listed
separately, and 6198 ordinary nonzero coordinate monomials after equal
coordinate monomials are collected. Exactly six nonconstant monomials
disappear. The constant coefficient decreases by three. Its unchanged
leading term is
\[
                    11082847123849826043664\cdot A^2.
\]
Exact integer comparison proves that the constant coefficient is
still the largest coefficient in absolute value and still has 1791357
binary digits. A retained coefficient equal to one gives content one;
a retained nonzero quartic monomial gives degree exactly four.

Run the following commands in the construction directory:
\begin{verbatim}
python3 build_quartic389.py
python3 check_quartic389.py
python3 verify_detector_normalizations.py
\end{verbatim}
The accompanying hashes bind the complete polynomial and quadratic
system. The finite computation checks the algebra and the asserted
numerical invariants. The local argument and its arithmetic inputs
are stated separately above.

\begin{thebibliography}{9}
\bibitem{Construction392}
J.~Brox and contributors, \emph{A forced prime condition and an explicit
quartic in 392 variables}, 8 November 2026, source and computational data
at commit \nolinkurl{de7aa60dd1794d7ea62a8a83cdc4bfd22380c74f},
\url{https://github.com/JoseBrox/diophantine-z-q_improvements}.
\bibitem{GlazerV13}
E.~Glazer and contributors, \emph{An explicit quartic defining the
integers in the rationals}, Version 13, commit
\nolinkurl{5bbd74f6a479776a332cbe036809d5dddae63720},
\url{https://github.com/elliotglazer/diophantine-z-q}.
\end{thebibliography}
\end{document}
